```latex
\chapter{Systems, Interfaces, and Feedback}
\label{chap:systems-interfaces-feedback}

\section{Why state changes category theory}
\label{sec:why-state-changes-category-theory}

Category theory is a theory of typed composition.  A morphism
\[
  f:X\to Y
\]
has a source, a target, and a rule for composing with other morphisms.  This
discipline is powerful because it separates the type of a process from the
process itself: a process can be composed precisely when its boundary matches
the boundary of another process.

Automata theory adds something that ordinary morphisms do not visibly carry:
state.  A stateful process does not merely transform an input boundary into an
output boundary.  It remembers, updates, responds, and exposes different
future behavior depending on its current state.  Once state is carried
categorically, composition is no longer only composition of arrows.  It is
composition of open systems whose internal states must be coupled along shared
interfaces.

This manuscript develops such a state-bearing categorical theory.  Its basic
morphisms are not ordinary arrows
\[
  A\to B,
\]
but Mealy morphisms
\[
  A\rightsquigarrow B.
\]
A Mealy morphism has an internal state object, typed input and output
interfaces, transition structure, and output structure.  It is therefore both
a categorical generalization of automata and a mathematical model of an open
system.

The central philosophical point of the manuscript is this:
\[
  \boxed{
  \text{category theory with persistent state naturally becomes
  categorical systems theory.}
  }
\]
This is not merely an analogy.  The formal objects developed below directly
realize the basic vocabulary of systems and cybernetics.  Interfaces are
internal categories.  Open systems are Mealy morphisms.  Maps of systems are
Mealy simulations.  Interconnection is pullback over shared interface data.
Behavior is residual semantics.  Behavioral predicates are subobjects of
behavior space, and derivative-closed predicates may be read as residual
behavioral contracts.  The theory also contains several distinct exact forms
of abstraction: representation as a Mealy morphism, semantic black-boxing,
Myhill--Nerode minimization, and contract-relative quotienting of operations.
Environments are downstream actions.  Feedback is trace over a hidden
interface.  Parallelism is product tensor.  True concurrency is cubical
execution.

Thus the abstract mathematics is not abstraction for its own sake.  It gives a
typed compositional language for systems with state, interaction, behavior,
feedback, and concurrent execution.

A first comparison is:
\[
\begin{array}{c|c|c}
\text{level}
&
\text{ordinary categorical form}
&
\text{state-bearing form in this manuscript}
\\
\hline
\text{interface}
&
\text{object}
&
\text{internal category}
\\
\text{process}
&
\text{morphism }X\to Y
&
\text{Mealy morphism }A\rightsquigarrow B
\\
\text{composition}
&
\text{ordinary composition}
&
\text{state coupling by pullback}
\\
\text{comparison}
&
\text{commuting diagram}
&
\text{Mealy simulation}
\\
\text{semantics}
&
\text{functorial interpretation}
&
\text{residual behavior assignment}
\\
\text{iteration}
&
\text{trace or fixed point}
&
\text{finite-path execution}
\\
\text{parallelism}
&
\text{monoidal tensor}
&
\text{product execution and cubical resolution}
\end{array}
\]

The conceptual movement can be displayed as:
\[
\begin{tikzcd}[column sep=huge,row sep=large]
  \text{ordinary morphism}
    \ar[r,"\text{add state}"]
  &
  \text{Mealy morphism}
    \ar[r,"\text{compose through interfaces}"]
  &
  \text{open system}
  \\
  X\to Y
    \ar[r,phantom,"\leadsto" description]
  &
  A\rightsquigarrow B
    \ar[r,phantom,"\leadsto" description]
  &
  \text{typed stateful interaction}
\end{tikzcd}
\]

The rest of the manuscript develops this idea formally.  This chapter explains
the systems-theoretic and cybernetic reading that motivates the construction.

\section{Systems and cybernetics as motivation}
\label{sec:systems-and-cybernetics-motivation}

Cybernetics began as the study of control and communication in animals and
machines.  Wiener's \emph{Cybernetics} gave the field its canonical
formulation as a study of messages, feedback, regulation, control, and
communication across biological and mechanical systems
\cite{Wiener1948Cybernetics}.  For the purposes of this manuscript, the
cybernetic lesson is that regulated behavior can often be analyzed through
loops of action, observation, correction, and further action.

Ashby gave cybernetics a sharper theory of state, transformation, black boxes,
regulation, and variety.  His \emph{Introduction to Cybernetics} develops a
formal vocabulary for adaptive and regulatory systems
\cite{Ashby1956IntroductionCybernetics}.  Ashby's law of requisite variety
says, roughly, that an effective regulator must possess enough variety to
counter the variety of relevant disturbances.  That law remains background
rather than a theorem of this manuscript.  Ashby's perspective is nevertheless
important: a system is partly characterized by the structured repertoire of
responses available to it.

General systems theory broadened the frame.  Bertalanffy argued that many
scientific objects should be understood as open systems: organized wholes
which persist through exchange with an environment
\cite{Bertalanffy1968GeneralSystemTheory}.  The relevant notions are system,
environment, boundary, exchange, organization, and emergence.  A living
organism, an ecological network, a firm, and a machine can share the same
structural problem: they maintain organization through interaction with what
lies outside them.

The behavioral approach to systems theory, developed by Willems and others,
moves still closer to the semantic side of this manuscript.  Willems argued
that systems should be described by their behavior: the set of admissible
trajectories or outcomes.  His mature account of open and interconnected
systems treats systems as things that are torn apart, linked, and reassembled
through shared variables \cite{Willems2007BehavioralApproach}.  His work on
interconnection by sharing variables is especially important for the present
viewpoint \cite{Willems2007InterconnectionSharingVariables}.  Interconnection
can mean identifying variables, imposing compatibility, or gluing systems
along a common boundary.

Adam and Dahleh extract an abstract categorical pattern from this behavioral
approach.  They emphasize that a behavioral system is naturally represented by
an inclusion
\[
  \mathcal B\hookrightarrow \mathbb U,
\]
where \(\mathbb U\) is a universe of possible outcomes and \(\mathcal B\) is
the admissible behavior.  They also distinguish syntax from semantics and
describe interconnection by variable sharing as a gluing operation
\cite{AdamDahleh2019AbstractBehavioralApproach}.  This is one of the closest
external precedents for the semantic parts of this manuscript.  Later chapters
treat specifications as subobjects of canonical behavior, systems as models
of those specifications, and interconnection as pullback-mediated
compatibility of typed interface data.

Categorical open-systems theory provides another point of contact.  Decorated
cospans, structured cospans, corelations, double categories, and hypergraph
categories have been used to model open networks and their compositional
semantics
\cite{Fong2015DecoratedCospans,Courser2020OpenSystemsDoubleCategories}.  This
literature uses technical foundations different from the span-based Mealy
calculus used here, while addressing the same philosophical problem: a theory
of systems must be a theory of composition through explicit boundaries.

Recent work in categorical systems theory makes this boundary-first
perspective explicit.  Myers emphasizes that a system interacts with its
environment through an interface that can be described separately from the
system itself \cite{MyersCategoricalSystemsTheory}.  Libkind and Myers
organize open systems, their interactions, and their maps over a double
category of interfaces and interactions
\cite{LibkindMyers2025DoubleSystems}.  The present manuscript uses a
different technical foundation: internal categories, spans, Mealy morphisms,
Mealy simulations, residual behavior, dynamic span logic, product
concurrency, polynomial response, and coproduct feedback.  The same
philosophical constraint is operative throughout: systems, interfaces,
interactions, and system maps are distinct layers of structure.

Categorical cybernetics is another recent development.  Work on lenses,
optics, open games, learning systems, and bidirectional processes studies
systems that interact with both environments and controllers
\cite{CapucciGavranovicHedgesRischel2021CategoricalCybernetics}.  The present
manuscript develops a span-based theory of typed Mealy systems.  Categorical
cybernetics confirms the same methodological lesson: cybernetic language has
mathematical force when attached to precise categorical structure.

\section{From automata to open systems}
\label{sec:from-automata-to-open-systems}

A classical automaton has states, inputs, transitions, and sometimes outputs.
A Mealy machine, in particular, updates its state while producing an output.
The theory developed here keeps that intuition but removes the restriction
that inputs and outputs are merely symbols.  Inputs and outputs are typed
events belonging to internal categories.

Thus an input interface is not just a set of letters.  It is an internal
category \(A\).  Its objects are boundary modes, and its arrows are admissible
input events between those modes.  Likewise, the output interface is an
internal category \(B\).

A system from \(A\) to \(B\) is then a Mealy morphism
\[
  \Phi:A\rightsquigarrow B.
\]
It has a carrier span
\[
\begin{tikzcd}[column sep=large]
  A_0
  &
  P
    \ar[l,"p"']
    \ar[r,"q"]
  &
  B_0 .
\end{tikzcd}
\]
The object \(P\) is the state object.  A state is not untyped: it has an
input-side exposure \(p(x)\in A_0\) and an output-side exposure
\(q(x)\in B_0\).

If \(a:u\to v\) is an input event of \(A\), and \(x\) is a state with
\(p(x)=v\), then the system may process \(a\).  It updates to a state
\[
  \delta(a,x)
\]
with input-side exposure
\[
  p(\delta(a,x))=u,
\]
and emits an output event
\[
  \omega(a,x):q(\delta(a,x))\to q(x)
\]
in \(B\).  This target-to-source orientation is used throughout the
manuscript.

A Mealy morphism is therefore a typed, categorical version of the
Mealy-machine idea.  But it is also more than automata theory.  Because
interfaces are categories and systems compose through spans, Mealy morphisms
form a calculus of open systems.

\section{Interfaces, systems, interactions, and maps}
\label{sec:interfaces-systems-interactions-maps}

The theory separates four levels that are often conflated in informal systems
language:
\[
  \text{interfaces,}
  \qquad
  \text{systems,}
  \qquad
  \text{interactions,}
  \qquad
  \text{maps of systems.}
\]
Each level has a distinct formal role.

\[
\begin{array}{c|c|c}
\text{layer}
&
\text{formal object}
&
\text{systems reading}
\\
\hline
\text{interface}
&
\text{internal category}
&
\text{typed boundary of interaction}
\\
\text{system}
&
\text{Mealy morphism}
&
\text{stateful open process inhabiting interfaces}
\\
\text{system map}
&
\text{Mealy simulation}
&
\text{structured comparison of systems}
\\
\text{interaction}
&
\text{composition, tensor, trace, execution}
&
\text{ways systems are coupled, run, hidden, or composed}
\\
\text{semantics}
&
\text{residual behavior}
&
\text{observable behavior under future interaction}
\end{array}
\]

Interfaces type interaction.  Systems inhabit interfaces.  Interactions
combine systems through interfaces.  Maps of systems compare systems while
respecting their transition-output structure.  Behavior records what remains
observable when the system is tested through future interface events.

This separation is the reason the formalism can be read systems-theoretically.
An open system is not merely a map.  It is a stateful process with a boundary.
Its boundary has its own transition structure.  Its internal state mediates
future response.  Its comparisons must preserve both update and output.  Its
behavior is obtained semantically rather than by inspecting implementation
witnesses.

\section{The systems dictionary}
\label{sec:systems-dictionary}

The principal systems-theoretic reading available for the formalism is
summarized by the following dictionary.  The entries concerning contracts and
abstraction are intended as translations into established systems language;
they do not require the technical development to replace its own terminology.
\[
\begin{array}{c|c}
\text{systems-theoretic language}
&
\text{categorical structure}
\\
\hline
\text{interface}
&
\text{internal category}
\\
\text{open system}
&
\text{Mealy morphism}
\\
\text{state}
&
\text{generalized element of the Mealy carrier}
\\
\text{input event}
&
\text{arrow of the input interface}
\\
\text{output event}
&
\text{arrow of the output interface}
\\
\text{interaction}
&
\text{operation combining systems through interfaces}
\\
\text{interconnection}
&
\text{pullback over shared interface data}
\\
\text{downstream environment}
&
\text{right action of the output interface}
\\
\text{one-step execution}
&
\text{primitive execution span}
\\
\text{execution structure}
&
\text{relative program category}
\\
\text{map of systems}
&
\text{Mealy simulation}
\\
\text{behavior space}
&
\text{canonical residual behavior object}
\\
\text{system behavior}
&
\text{semantic image of the behavior assignment}
\\
\text{state observation}
&
\text{behavior assignment}
\\
\text{behavioral predicate or raw contract}
&
\text{subobject of canonical behavior}
\\
\text{residual behavioral contract}
&
\text{derivative-closed coequational subobject}
\\
\text{contract presentation}
&
\text{map into canonical behavior}
\\
\text{contract satisfaction}
&
\text{factorization of the semantic map}
\\
\text{structural abstraction}
&
\text{representation as a Mealy morphism}
\\
\text{behavioral or black-box abstraction}
&
\text{semantic image of a system}
\\
\text{minimal residual abstraction}
&
\text{Myhill--Nerode quotient and minimal realization}
\\
\text{contract-relative operational abstraction}
&
\text{quotient by the operation congruence}
\\
\text{may modality}
&
\text{diamond}
\\
\text{must modality}
&
\text{box}
\\
\text{guard}
&
\text{subobject represented as a test span}
\\
\text{iteration}
&
\text{finite-path construction}
\\
\text{feedback}
&
\text{coproduct trace over a hidden interface}
\\
\text{choice}
&
\text{homwise coproduct}
\\
\text{parallelism}
&
\text{product tensor}
\\
\text{true concurrency}
&
\text{cubical resolution of product execution}
\\
\text{response repertoire}
&
\text{polynomial response profile}.
\end{array}
\]

The table should be read as a disciplined interpretation rather than as a
loose metaphor.  The point is not that every word from systems theory is
automatically available.  The point is that the central structural words of
systems theory---interface, system, interaction, behavior, specification,
contract, abstraction, environment, feedback, and concurrency---have precise
translations in the categorical automata developed here.  In particular,
``contract'' and ``abstraction'' are not single additional structures imposed
on the theory.  They name several structures that are already distinguished
categorically.

\section{Composition as interconnection}
\label{sec:composition-as-interconnection}

The basic operation on open systems is composition.  Suppose
\[
  P:A\rightsquigarrow B,
  \qquad
  Q:B\rightsquigarrow C.
\]
The two systems share the interface \(B\).  The output boundary of \(P\) must
match the input boundary of \(Q\).  Categorically, this matching is expressed
by the pullback
\[
\begin{tikzcd}[column sep=large,row sep=large]
  P\times_{B_0}Q
    \ar[r]
    \ar[d]
  &
  Q
    \ar[d,"p_Q"]
  \\
  P
    \ar[r,"q_P"']
  &
  B_0 .
\end{tikzcd}
\]
Thus a state of the composite is a compatible pair \((x,y)\) with
\[
  q_P(x)=p_Q(y).
\]

This is the formal meaning of interconnection in the span-based Mealy
calculus:
\[
  \boxed{
  \text{interconnection is pullback over shared interface data.}
  }
\]
Systems do not compose by forgetting their boundaries.  They compose by
identifying the boundary data through which one system can drive, constrain,
or inform another.

This gives a categorical version of a central systems-theoretic idea.
Willems's behavioral approach treats interconnection as sharing variables
\cite{Willems2007InterconnectionSharingVariables}.  In the present manuscript,
shared variables are replaced by typed interface data, and compatibility is
expressed by pullback.

\section{Environments and execution}
\label{sec:environments-and-execution}

An open system has an output interface.  An environment is something on which
outputs can act.

If
\[
  \Phi:A\rightsquigarrow B
\]
is an open system, then \(B\) is the output interface.  A downstream
environment for \(\Phi\) is represented by a right internal \(B\)-action
\(Y\).  The coupled state object is
\[
  S_{\Phi,Y}=P\times_{B_0}Y.
\]
A point of \(S_{\Phi,Y}\) is a system state together with a compatible
environment state:
\[
\begin{tikzcd}[column sep=large,row sep=large]
  S_{\Phi,Y}
    \ar[r]
    \ar[d]
  &
  Y
    \ar[d]
  \\
  P
    \ar[r,"q"']
  &
  B_0 .
\end{tikzcd}
\]

When the system emits a \(B\)-arrow, that arrow acts on the environment
component.  The result is a right \(A\)-action, whose action category is the
relative program category
\[
  \mathsf{Prog}_{\Phi,Y}.
\]

Execution therefore belongs to the system-in-environment.  The Mealy morphism
is the open system.  The pair \((\Phi,Y)\) is the system placed downstream in
an environment.  The relative program category is its generated execution
structure.

This gives a precise cybernetic reading of output.  An output is not merely a
label.  It is an interface event capable of acting on a downstream system.

\section{Maps and simulations}
\label{sec:maps-and-simulations}

A theory of systems needs maps of systems.  In this manuscript, maps of
systems are Mealy simulations.

A displayed simulation
\[
  \Theta:\Phi\Rightarrow\Psi
\]
is a structured comparison of Mealy morphisms.  Its variance is important.
Although the displayed simulation points from \(\Phi\) to \(\Psi\), its
carrier comparison is directed from the carrier of \(\Psi\) to the carrier of
\(\Phi\), together with an output-comparison arrow in the output interface.

This variance is not an accident.  It is what makes simulations act
contravariantly on evaluated execution.  For a fixed downstream action \(Y\),
a simulation
\[
  \Theta:\Phi\Rightarrow\Psi
\]
induces a functor
\[
  H_{\Theta,Y}:
  \mathsf{Prog}_{\Psi,Y}
  \longrightarrow
  \mathsf{Prog}_{\Phi,Y}.
\]
Thus comparison of systems transports executions in the opposite direction.

This is the comparison discipline used throughout the manuscript: displayed
simulation direction and execution-transport direction are opposite at the
level of evaluated program categories.  A Mealy simulation is therefore not
just a map of state sets.  It is a structured comparison respecting input
processing and output comparison.

\section{Behavior and black-boxing}
\label{sec:behavior-and-black-boxing}

A system has internal state, but its externally meaningful content is its
behavior.  In this manuscript, behavior is residual: it records what the
system can do under all future typed interactions.

For fixed interfaces \(A\) and \(B\), there is a canonical behavior object
\[
  \mathrm{Beh}^{\mathrm{can}}_{A,B}.
\]
Every Mealy morphism
\[
  P:A\rightsquigarrow B
\]
has a semantic map
\[
  \beta_P:P\longrightarrow \mathrm{Beh}^{\mathrm{can}}_{A,B}.
\]
This map sends a state to its residual behavior.

The semantic image
\[
  \mathrm{Beh}(P)\hookrightarrow \mathrm{Beh}^{\mathrm{can}}_{A,B}
\]
is the behavior generated by the system.  Passing from \(P\) to
\(\mathrm{Beh}(P)\) is a categorical black-boxing operation.  The internal
implementation witnesses in \(P\) are forgotten; the residual observable
behavior remains.

The same theme appears in the Myhill--Nerode construction.  States that
determine the same residual behavior are identified.  The resulting minimal
residual realization is the behaviorally separated system obtained by
quotienting away redundant implementation state.

This is where the present theory meets the behavioral tradition in systems
theory most directly.  A system is not only a state machine.  It is also a
generator of admissible behavior inside a universe of possible behavior
\cite{Willems2007BehavioralApproach,AdamDahleh2019AbstractBehavioralApproach}.

\section{Contracts and the several forms of abstraction}
\label{sec:contracts-and-abstraction}

The words ``contract'' and ``abstraction'' are used in several neighboring
areas of systems theory, but neither word has a unique mathematical meaning.
The constructions of this manuscript admit precise translations into that
language.  This section records those translations without changing the
technical names of the constructions themselves.

\subsection{Behavioral predicates and residual contracts}
\label{subsec:behavioral-predicates-residual-contracts}

Fix interfaces \(A\) and \(B\), put
\[
  Z=A_0\times B_0,
\]
and write
\[
  \mathrm{Beh}_0
\]
for the carrier of the canonical behavior object
\[
  \mathrm{Beh}^{\mathrm{can}}_{A,B}.
\]
An arbitrary subobject
\[
  C\hookrightarrow \mathrm{Beh}_0
\]
in \(\mathcal E/Z\) is a \emph{behavioral predicate}: it selects a class of
admissible residual behaviors.  In systems language, such a predicate may be
read as a \emph{raw behavioral contract}.  It states which behaviors are
allowed, but it need not yet be stable under continued interaction.

A behavioral predicate
\[
  V\hookrightarrow \mathrm{Beh}_0
\]
is \emph{residually invariant} when it is closed under every canonical
derivative:
\[
  H\in V_v,\quad a:u\to v
  \quad\Longrightarrow\quad
  \partial_aH\in V_u.
\]
Equivalently, the canonical derivative map restricts to \(V\).  The technical
term for such a predicate is a \emph{local coequational theory}.  Its
systems-theoretic translation is a \emph{residual behavioral contract}.  The
word ``residual'' records the decisive closure condition: a behavior admitted
now must continue to satisfy the contract after every typed input event for
which a residual behavior is formed.

A map
\[
  k:K\to \mathrm{Beh}_0
\]
is more general than a subobject.  It presents a family of designated
behaviors and may therefore be read as a \emph{contract presentation}.  Its
derivative image
\[
  \mathrm{Der}_0(K)
\]
is the least residual behavioral contract containing the image of \(k\).  In
this interpretation, \(\mathrm{Der}_0(K)\) is the \emph{generated residual
contract} of the presentation.

There are in fact two canonical ways to pass between arbitrary behavioral
predicates and residual contracts.  Let
\[
  \mathsf{Pred}_{A,B}
  =
  \operatorname{Sub}_{\mathcal E/Z}(\mathrm{Beh}_0)
\]
be the poset of behavioral predicates, and let
\[
  \mathsf{RCtr}_{A,B}
  \subseteq
  \mathsf{Pred}_{A,B}
\]
be the full subposet of residually invariant predicates.  Write
\[
  I_{\partial}:
  \mathsf{RCtr}_{A,B}
  \hookrightarrow
  \mathsf{Pred}_{A,B}
\]
for the inclusion.

For a predicate \(C\), its \emph{generated residual closure}
\[
  \operatorname{Gen}_{\partial}(C)
\]
is the least residually invariant predicate containing \(C\).  It is computed
by the derivative image:
\[
  \operatorname{Gen}_{\partial}(C)
  =
  \operatorname{im}
  \left(
    A_1\times_{t_A,A_0,p_C}C
    \xrightarrow{\;\partial(1\times j_C)\;}
    \mathrm{Beh}_0
  \right).
\]
Thus \(\operatorname{Gen}_{\partial}\) freely adds every residual consequence
forced by the behaviors already admitted by \(C\).

Dually, let
\[
  D\mathrm{Beh}
  =
  A_1\times_{t_A,A_0,p_{\mathrm{Beh}}}\mathrm{Beh}_0,
\]
and write
\[
  \pi_{\mathrm{Beh}}:D\mathrm{Beh}\to\mathrm{Beh}_0,
  \qquad
  \partial:D\mathrm{Beh}\to\mathrm{Beh}_0
\]
for the projection to the original behavior and the canonical derivative.
The \emph{residual safety core} of \(C\) is
\[
  \operatorname{Safe}_{\partial}(C)
  =
  \forall_{\pi_{\mathrm{Beh}}}\bigl(\partial^\ast C\bigr).
\]
Here the pullback and dependent universal quantification are taken in
\(\mathcal E\), and the resulting subobject of \(\mathrm{Beh}_0\) is viewed
over \(Z\).  Internally,
\[
  H\in\operatorname{Safe}_{\partial}(C)
  \quad\Longleftrightarrow\quad
  \forall a:u\to p_{\mathrm{Beh}}(H),\;
  \partial_aH\in C.
\]
It is the greatest residually invariant predicate contained in \(C\).  It
removes precisely those behaviors having some residual continuation that
violates the original predicate.

These constructions form the adjoint triple
\[
  \boxed{
  \operatorname{Gen}_{\partial}
  \dashv
  I_{\partial}
  \dashv
  \operatorname{Safe}_{\partial}.
  }
\]
Consequently,
\[
  \operatorname{Safe}_{\partial}(C)
  \leq
  C
  \leq
  \operatorname{Gen}_{\partial}(C),
\]
and the residual behavioral contracts are exactly the common fixed points:
\[
  \operatorname{Gen}_{\partial}(V)=V
  \quad\Longleftrightarrow\quad
  V\text{ is residually invariant}
  \quad\Longleftrightarrow\quad
  \operatorname{Safe}_{\partial}(V)=V.
\]
The left adjoint gives the least contract obtained by closing a predicate
under residual behavior.  The right adjoint gives the largest contract that
can safely be asserted inside a possibly non-invariant predicate.

When the ambient subobject lattices are Heyting algebras, an
assumption--guarantee pair can also be represented at the predicate level.
For behavioral predicates
\[
  \mathsf{Assm},\mathsf{Guar}
  \hookrightarrow
  \mathrm{Beh}_0,
\]
the implication
\[
  \mathsf{Assm}\Rightarrow\mathsf{Guar}
\]
is a raw assume--guarantee contract, while
\[
  \operatorname{Safe}_{\partial}
  \bigl(\mathsf{Assm}\Rightarrow\mathsf{Guar}\bigr)
\]
is its largest residual behavioral contract.  This translation locates
assumption--guarantee reasoning inside the semantic doctrine; it does not by
itself assert a complete assume--guarantee calculus.

A system
\[
  P:A\rightsquigarrow B
\]
satisfies a residual behavioral contract \(V\) precisely when its semantic map
factors through \(V\):
\[
\begin{tikzcd}[column sep=large,row sep=large]
  P
    \ar[rr,"\beta_P"]
    \ar[dr,dashed,"\bar\beta_P"']
  &&
  \mathrm{Beh}^{\mathrm{can}}_{A,B}
  \\
  &
  V
    \ar[ur,hook,"j_V"']
  &
\end{tikzcd}
\]
Equivalently,
\[
  \mathrm{Beh}(P)\leq V.
\]
Thus the usual model relation may be read as \emph{contract satisfaction}.
The semantic image \(\mathrm{Beh}(P)\) is the least residual behavioral
contract through which the semantics of \(P\) factors, and may therefore be
read as the \emph{principal residual contract generated by the system}.
Conversely, the restriction of the canonical behavior machine to \(V\) is the
canonical realization of that contract.

\subsection{Several exact meanings of abstraction}
\label{subsec:several-exact-meanings-abstraction}

The theory contains several mathematically distinct operations that may all be
described informally as abstraction.  They should not be conflated.

\paragraph{Structural abstraction.}
A concrete stateful mechanism is structurally abstracted when it is represented
as a Mealy morphism
\[
  P:A\rightsquigarrow B.
\]
This forgets implementation detail not represented by the chosen state object,
interfaces, transition, and output maps.  It is abstraction into the formal
language of the theory rather than a quotient constructed inside that
language.

\paragraph{Behavioral or black-box abstraction.}
The terminal semantic map has a regular epi--mono factorization
\[
\begin{tikzcd}[column sep=large,row sep=large]
  P
    \ar[rr,"\beta_P"]
    \ar[dr,two heads,"\eta_P"']
  &&
  \mathrm{Beh}^{\mathrm{can}}_{A,B}
  \\
  &
  \mathrm{Beh}(P)
    \ar[ur,hook]
  &
\end{tikzcd}
\]
The passage
\[
  P\twoheadrightarrow\mathrm{Beh}(P)
\]
forgets distinctions between states that have the same residual semantics.
This is the basic \emph{behavioral abstraction} or \emph{black-box
abstraction} of a system.

\paragraph{Minimal residual abstraction.}
For a behavior or contract presentation
\[
  k:K\to\mathrm{Beh}_0,
\]
the formal derivative orbit \(D_K\) carries a residual evaluation map whose
kernel pair is the Myhill--Nerode equivalence
\[
  {\equiv_K}.
\]
The quotient satisfies
\[
  D_K/{\equiv_K}
  \cong
  \mathrm{Der}_0(K),
\]
and the corresponding restricted canonical machine
\[
  \mathrm{Min}(K)
\]
is the minimal reachable behaviorally separated realization of the generated
residual contract.  This is a stronger abstraction than merely taking the
semantic image of a pre-existing implementation: it constructs the canonical
minimal residual realization of the presented behavior.

\paragraph{Contract-relative operational abstraction.}
A residual behavioral contract
\[
  V\hookrightarrow\mathrm{Beh}_0
\]
determines an observational congruence \(R_V\) on input operations.  Two
suitably typed operations are identified when they induce the same
transition--output action on every behavior admitted by \(V\):
\[
  a\,R_V\,a'
  \quad\Longleftrightarrow\quad
  \forall H\in V,\;
  \chi_V(a,H)=\chi_V(a',H).
\]
The quotient
\[
  \mathrm{Op}(V)
  \cong
  A^{\mathrm{op}}/R_V
\]
is therefore the \emph{contract-relative operational abstraction}: it retains
exactly the operational distinctions observable by behaviors satisfying
\(V\).  For a presentation \(k\), the syntactic transition--output category
\[
  \mathrm{SynTO}(K)
  =
  \mathrm{Op}\bigl(\mathrm{Der}_0(K)\bigr)
\]
is the corresponding syntactic abstraction of the generated residual
contract.

These meanings of abstraction act on different data:
\[
\begin{array}{c|c|c}
\text{form of abstraction}
&
\text{data being abstracted}
&
\text{result}
\\
\hline
\text{structural}
&
\text{concrete mechanism}
&
\text{Mealy morphism}
\\
\text{behavioral or black-box}
&
\text{implementation states}
&
\mathrm{Beh}(P)
\\
\text{minimal residual}
&
\text{formal residual presentations}
&
\mathrm{Min}(K)
\\
\text{contract-relative operational}
&
\text{input operations}
&
\mathrm{Op}(V)
\end{array}
\]
They are exact images or quotients.  They are not, merely by being called
abstractions, instances of approximate abstract interpretation.  Such an
interpretation would require additional concrete and abstract domains, an
approximation order, and a soundness relation.

The technical vocabulary of residual behavior, semantic image,
Myhill--Nerode quotient, local coequational theory, operation congruence, and
syntactic transition--output category remains more discriminating than the
single words ``contract'' and ``abstraction.''  The purpose of the present
translation is therefore interpretive: it identifies how established systems
language maps onto the categorical structures without requiring those
structures to be renamed elsewhere.

\section{Specifications, contracts, and requirements}
\label{sec:specifications-and-requirements}

Specifications are behavioral constraints.  An arbitrary specification may be
represented by a behavioral predicate
\[
  C\hookrightarrow \mathrm{Beh}_0.
\]
In the translation established above, this is a raw behavioral contract.  A
specification stable under every residual derivative is a local coequational
theory
\[
  V\hookrightarrow \mathrm{Beh}^{\mathrm{can}}_{A,B},
\]
and may be read as a residual behavioral contract.

A system
\[
  P:A\rightsquigarrow B
\]
satisfies the specification when its semantic map factors through \(V\):
\[
\begin{tikzcd}[column sep=large,row sep=large]
  P
    \ar[rr,"\beta_P"]
    \ar[dr,dashed]
  &&
  \mathrm{Beh}^{\mathrm{can}}_{A,B}
  \\
  &
  V
    \ar[ur,hook]
  &
\end{tikzcd}
\]
The class of all such systems is the model class of \(V\).  In contract
language, the same factorization says that \(P\) satisfies \(V\).  The
terminological translation adds no new satisfaction relation: modelhood and
contract satisfaction are the same semantic factorization.

The distinction between a raw predicate and a residual contract clarifies two
ways of repairing a specification that is not derivative-closed.  Its
generated residual closure
\[
  \operatorname{Gen}_{\partial}(C)
\]
adds all residual consequences and is the least residual contract containing
\(C\).  Its residual safety core
\[
  \operatorname{Safe}_{\partial}(C)
\]
removes unsafe behaviors and is the greatest residual contract contained in
\(C\).  These are respectively the permissive closure and conservative core
of the original requirement.

Requirements are semantic before they are syntactic.  Syntax is extracted
from behavior through transition-output images, operational congruences, and
syntactic transition-output categories.  This reverses a common order of
presentation.  Rather than beginning with an algebra of syntax and then
assigning semantics, the manuscript first constructs behavior space and then
extracts the algebraic syntax recognized by behavioral constraints.

The Birkhoff theorems express this idea categorically.  Behaviorally
definable classes of systems are exactly those closed under the appropriate
system-theoretic operations: homomorphic preimage, quotient, coproduct, and
reindexing, depending on the local or global setting.  Under the contract
translation, these are contract-definability theorems: they characterize the
classes of systems definable by residual behavioral contracts.  Under the
abstraction translation, the semantic image supplying the defining theory is
the behavioral abstraction generated by the class.  These readings are
interpretive descriptions of the same categorical results, not alternative
technical formulations.

\section{Observation and dynamic logic}
\label{sec:observation-and-dynamic-logic}

Observation appears in several related forms.  At the semantic level,
observation is the behavior assignment.  At the execution level, observation
is given by output labelling and downstream comparison.  At the logical level,
observation is represented by predicate transformers.

For a span
\[
  R=
  \left(
    X\xleftarrow{\ell}R\xrightarrow{r}Y
  \right),
\]
the diamond modality is
\[
  \langle R\rangle\varphi
  :=
  \exists_{\ell}(r^\ast\varphi),
\]
and the box modality is
\[
  [R]\varphi
  :=
  \forall_{\ell}(r^\ast\varphi).
\]
The diamond is the may modality: there exists an \(R\)-execution reaching
\(\varphi\).  The box is the must modality: every \(R\)-execution satisfies
\(\varphi\).

The modal reading is:
\[
\begin{array}{c|c|c}
\text{logical form}
&
\text{categorical definition}
&
\text{systems reading}
\\
\hline
\langle R\rangle\varphi
&
\exists_{\ell}(r^\ast\varphi)
&
\text{some execution reaches }\varphi
\\
\left[R\right]\varphi
&
\forall_{\ell}(r^\ast\varphi)
&
\text{all executions satisfy }\varphi
\\
T_\theta
&
S\xleftarrow{\theta}\Theta\xrightarrow{\theta}S
&
\text{test or guard}
\\
\mathsf{Path}(U)
&
\coprod_{n\geq 0}U^{[n]}
&
\text{finite iteration}
\end{array}
\]

A test is a predicate regarded as a program.  If
\[
  \theta:\Theta\hookrightarrow S
\]
is a subobject, its test span is
\[
  T_\theta=
  \left(
    S\xleftarrow{\theta}\Theta\xrightarrow{\theta}S
  \right).
\]
Finite iteration is represented by the finite-path construction
\[
  \mathsf{Path}(U)
  =
  \coprod_{n\geq 0}U^{[n]}.
\]
Its modal semantics is fixed-point semantics:
\[
  \langle\mathsf{Path}(U)\rangle\varphi
  =
  \mu X.(\varphi\vee\langle U\rangle X),
\]
and
\[
  [\mathsf{Path}(U)]\varphi
  =
  \nu X.(\varphi\wedge[U]X).
\]
Thus may, must, test, guard, and finite iteration have direct categorical
content.

\section{Feedback as hidden-interface closure}
\label{sec:feedback-as-hidden-interface-closure}

Feedback is the cybernetic idea most directly realized by the mathematics of
this manuscript.  A feedback situation separates visible interface from
hidden interface:
\[
  X+H\longrightarrow Y+H.
\]
The visible source is \(X\), the visible target is \(Y\), and \(H\) is the
interface through which execution may circulate internally.

The block form is
\[
  R=
  \begin{pmatrix}
    R_{XY} & R_{XH}\\
    R_{HY} & R_{HH}
  \end{pmatrix}.
\]
The block \(R_{XH}\) enters the hidden interface, \(R_{HH}\) moves within the
hidden interface, \(R_{HY}\) exits the hidden interface, and \(R_{XY}\) exits
immediately through the visible target.

The feedback trace over \(H\) is
\[
  \operatorname{Tr}^{H}(R)
  =
  R_{XY}
  \boxplus
  R_{HY}\circ
  \mathsf{Path}(R_{HH})
  \circ
  R_{XH}.
\]
Thus a visible execution either exits immediately through \(R_{XY}\), or it
enters the hidden interface, follows finitely many hidden \(H\)-steps, and
then exits through \(Y\).

\[
  \boxed{
  \text{feedback is finite hidden-path closure.}
  }
\]

This is why the trace construction is not merely a formal categorical
operation.  It is the mathematical form of closing an internal loop while
retaining the externally visible behavior of the resulting system.  The
formal trace used later is the coproduct trace on spans, related to the
general theory of traced monoidal categories
\cite{JoyalStreetVerity1996TracedMonoidalCategories} and to categorical
feedback semantics
\cite{KatisSabadiniWalters2002FeedbackTraceFixedPoint}.

Guarded feedback is obtained by inserting tests.  A guarded while program has
the form
\[
  (\theta?;u)^{\star_{\mathrm{path}}};\epsilon?,
\]
where \(\theta\) is the continuation guard and \(\epsilon\) is the exit
condition.  Its diamond semantics is a least fixed point, and its box
semantics is a greatest fixed point.  This is the point at which cybernetic
feedback, program iteration, and dynamic logic meet.

\section{Choice, parallelism, and true concurrency}
\label{sec:choice-parallelism-true-concurrency}

The theory distinguishes choice, parallelism, and true concurrency.

Choice is homwise coproduct.  If
\[
  R,S:X\to Y
\]
are parallel spans, then
\[
  R\boxplus S:X\to Y
\]
is the span whose witnesses are either \(R\)-witnesses or \(S\)-witnesses.
This is a witness-retaining form of nondeterministic alternative.

Parallelism is product tensor.  Products expose simultaneous coordinates.  The
product tensor places systems side by side.

True concurrency is obtained by resolving product execution into cubical
execution.  A product of independent systems gives independent execution
coordinates.  The cubical construction records filled cells together with
their boundary paths.  Boundary linearizations are sequential observations of
a higher-dimensional execution.

The distinction is:
\[
\begin{array}{c|c|c}
\text{systems notion}
&
\text{categorical operation}
&
\text{interpretation}
\\
\hline
\text{choice}
&
R\boxplus S
&
\text{alternative witnesses}
\\
\text{parallelism}
&
R\times S
&
\text{simultaneous coordinates}
\\
\text{true concurrency}
&
\text{cubical execution}
&
\text{higher-dimensional event structure}
\\
\text{feedback}
&
\operatorname{Tr}^{H}(R)
&
\text{hidden-interface closure}
\end{array}
\]

This distinction keeps concurrency separate from nondeterministic
interleaving.  Parallel systems do not merely interleave their steps.  Their
independent coordinates support higher-dimensional cells, and the cubical
semantics records that independence.

\section{Response repertoires}
\label{sec:response-repertoires}

The response-profile chapter supplies one more cybernetic term with a precise
formal reading:
\[
  \text{response repertoire}
  =
  \text{one-step polynomial response profile}.
\]
A response profile records the one-step query-response structure available at
a state.  It is obtained by currying relative Mealy execution data.  Flattening
the response profile returns the primitive execution span.

The word ``repertoire'' is appropriate here.  The response profile records the
structured set of possible one-step responses available to the system at a
state and interface.  It is related in spirit to cybernetic discussions of
variety, while the stronger phrase ``requisite variety'' is reserved for a
setting with an additional theorem comparing disturbance variety with response
variety.

\section{What the theory does not yet claim}
\label{sec:what-the-theory-does-not-yet-claim}

The systems reading is precise, but it is not unlimited.  The manuscript
formalizes typed openness, interaction, behavior, specification, execution,
feedback, and concurrency.  It does not automatically formalize every term
from cybernetics or systems theory.

For example, homeostasis would require a theorem saying that some controlled
quantity is maintained under a class of disturbances.  Learning would require
training, update, or adaptation dynamics.  Agency would require goals,
preferences, choices, or evaluation structure.  Autopoiesis would require a
theory of self-production.  Ashby's law of requisite variety would require a
formal comparison between disturbance variety and response variety.

These are not rejected.  They are simply not smuggled in by terminology.  The
method of the manuscript is to use systems-theoretic language exactly where
the categorical structure supports it.

\section{Reading guide}
\label{sec:systems-reading-guide}

The chapters that follow develop the formal machinery behind the dictionary.
The contract and abstraction vocabulary recorded in
Section~\ref{sec:contracts-and-abstraction} is a translation of that machinery;
the technical development may retain the more specific terms appropriate to
each construction.

The first technical layer constructs Mealy morphisms between internal
categories and Mealy simulations between them.  This gives the basic calculus
of typed open systems and their maps.

The second layer develops residual behavior, canonical behavior objects,
semantic maps, Myhill--Nerode quotients, and minimal residual realizations.
This is the semantic and behavioral part of the theory.

The third layer studies specifications as coequational subobjects of behavior
and proves Birkhoff-style theorems relating behavioral theories to closure
properties of system classes.

The fourth layer extracts algebraic syntax from behavior through
transition-output images, operational congruences, and syntactic
transition-output categories.

The fifth layer places systems in downstream environments, producing relative
program categories and dynamic span logic.

The sixth layer curries relative execution into polynomial response profiles,
making precise the one-step response repertoire of a system.

The seventh layer studies product execution and cubical concurrency,
distinguishing true concurrent structure from mere interleaving.

The feedback chapter then studies coproduct trace, finite hidden paths,
guarded iteration, signed interfaces, and feedback semantics.

Together these layers support the thesis of this chapter:
\[
  \boxed{
  \text{stateful categorical automata form a typed categorical theory of
  open systems.}
  }
\]

\section{Summary}
\label{sec:systems-summary}

This manuscript develops a typed open-systems calculus.

Internal categories are interfaces.  Mealy morphisms are open systems
inhabiting pairs of interfaces.  Mealy simulations are maps or comparisons of
systems.  Interactions are the operations that compose, couple, tensor,
execute, iterate, or trace systems.  Pullbacks express interconnection by
shared interface data.  Downstream actions provide environments.  Relative
program categories describe execution in an environment.  Residual semantics
gives behavior.  Behavioral predicates give raw contracts, while
derivative-closed coequational subobjects give residual contracts and
specifications.  Semantic images, Myhill--Nerode quotients, and operational
congruences provide distinct exact forms of abstraction.  Homwise coproduct
gives choice.  Product tensor gives parallelism and, after cubical resolution,
true concurrency.  Coproduct trace gives feedback by finite hidden-path
closure.

The point is not metaphorical.  The categorical structures themselves carry
the systems-theoretic interpretation.  Category theory with persistent state
naturally becomes a theory of typed open systems, their interfaces, their
interactions, their maps, their behavior, their feedback, and their concurrent
execution.
```
